\begin{textsample}{POS Dim 1 – vem\_ed\_02 – Score 2.00 – vem\_ed\_02/t006.txt}  \label{ex:f1_pos_163}
EM BUSCA DA INTERNACIONALIZAÇÃO Em 01 de julho de 2020, o professor Michael Fields, da Eastern Oregon University (EUA), apresentou uma videoconferência para as Fatecs Ferraz de Vasconcelos, Itaquaquecetuba e Mogi das Cruzes.
Com o \textbf{apoio} de Geraldo Ribeiro Filho, \textbf{que} leciona nessas três Fatecs, o evento \textbf{foi} organizado pela equipe dos Projetos Colaborativos Internacionais (PCI – Cesu), por meio de seu coordenador, Osvaldo Succi Junior, e da professora Neusa Haruka Sezaki Gritti.
Essa iniciativa busca incentivar a \textbf{integração} das unidades das Fatecs no \textbf{processo} de internacionalização de seus currículos e campi.
Doutor em Global Leadership, coach de esportes coletivos e pai de cinco filhos entre 10 e 18 anos, Michael Fields nasceu em Indianápolis, estado de Indiana; viveu dois anos no Novo México e atualmente mora e trabalha em La Grande, cidade de 15 mil habitantes no estado do Oregon (noroeste dos Estados Unidos).
Apresentou a diversidade geográfica, demográfica e econômica dessas três regiões dos Estados Unidos, e depois falou sobre a Eastern Oregon University, os cursos \textbf{que} \textbf{são} oferecidos e o calendário acadêmico.
Assistiram a palestra 40 \textbf{pessoas}, entre elas o diretor da Fatec Mogi das Cruzes, Bruno Marques Panccioni; os coordenadores da Fatec Itaquaquecetuba Francisco Claudio Tavares (Gestão Comercial) e Aparecido Rodrigues da Silva López-Guerrero (Gestão da Tecnologia da Informação); a coordenadora da Fatec Ferraz de Vasconcelos Andrea Zotovici (Análise e Desenvolvimento de Sistemas), além de professores e alunos das unidades.
O professor Wilton Garcia, da Fatec Itaquaquecetuba, acompanhou a videoconferência e divulgou em https://fatecitaqua.wordpress.com/ QUEM \textbf{É} QUEM Sally Crimmins-Villela, vice-chanceler associada para assuntos globais e executiva internacional sênior da State University of New York (SUNY), supervisiona \textbf{programas} acadêmicos internacionais, internacionalização \textbf{curricular} e “Learning Through Development”, um paradigma por meio do qual universidades apoiam o \textbf{desenvolvimento} voltado à comunidade.
Orienta operações da SUNY no México, Turquia e Haiti, e conduz o SUNY COIL Center, \textbf{referência} mundial em \textbf{intercâmbios} virtuais, infundindo aprendizagem global em 18 países e mais de 70 instituições (dentre as quais estão as Fatecs).
Falante fluente de português e espanhol, Sally estudou, trabalhou e viveu no Brasil, no México, na Espanha e em Portugal.
Lecionou Língua Portuguesa no Union College e na SUNY at Albany (ambas instituições, no Estado de Nova York).
Em sua temporada brasileira, ensinou Língua Inglesa na Universidade Federal do Rio Grande do Norte (UFRN).
Ao receber a primeira edição do VEm com PCI, comentou: “As Fatecs \textbf{são} líderes em Virtual Exchange (\textbf{intercâmbios} virtuais) não só no Brasil, mas no mundo, \textbf{são} uma inspiração para \textbf{todos} nós”.

% matched lemmas: apoio, curricular, desenvolvimento, integração, intercâmbio, pessoa, processo, programa, que, referência, ser, todo
\end{textsample}
